\begin{center}
    {\large \textbf{The Epidemiological Cost of Behavioral Delay: Integrating Intertemporal Preferences into SEIR Dynamics for Vaccination Policy}}
\end{center}

\vspace{1em}

\begin{abstract}
\noindent \textbf{Background:} While vaccination coverage is a standard metric for public health success, the \textit{timing} of uptake---often hindered by administrative and behavioral delays---critically determines outbreak trajectories. Traditional epidemiological models often treat vaccination rates as supply-driven constants, neglecting the structural time preferences of the population.

\noindent \textbf{Methods:} We develop a behaviorally-augmented SEIR model that endogenizes vaccination uptake as a function of individual delay aversion ($\kappa$). We recover structural discounting parameters from a large-scale discrete choice experiment ($N=1,027$) conducted in Wuhan, China. These parameters are integrated into the transmission dynamics to quantify the ``behavioral gap'' between supply capacity and actual uptake.

\noindent \textbf{Results:} Our empirical results identify significant present-biased preferences ($\hat{\kappa}=0.225$) in the population. Simulation results demonstrate that this behavioral delay aversion can lead to a ``behavioral surge'' in infections: compared to a supply-only constrained scenario (Attack Rate $\approx$ 22.5\%), the introduction of behavioral delay leads to a dramatic expansion of the outbreak (Attack Rate $\approx$ 63.6\%), representing a 183\% increase in cumulative infections. Subgroup analysis reveals that older adults ($\hat{\kappa}=0.327$) and those with neutral institutional trust ($\hat{\kappa}=0.356$) exhibit the highest behavioral hurdles, translating to the greatest biological vulnerability.

\noindent \textbf{Conclusion:} Behavioral discounting is a structural driver of epidemic risk. Public health interventions that solely expand supply without addressing the ``psychological cost of waiting'' are likely to underperform. Preventive strategies must prioritize behavioral ``de-biasing'' and administrative streamlining to bridge the gap between vaccine availability and timely biological protection.

\vspace{1em}
\noindent \textbf{Keywords:} Vaccination timing; Behavioral SEIR; Intertemporal choice; Delay aversion; Outbreak dynamics; Public health policy.
\end{abstract}
